% Section 1: Introduction

\subsection{The Execution-Layer Fallacy}

The dominant paradigm for deploying AI systems treats them as execution layers: a user provides a task, the AI executes it, and returns output. This model---borrowed from traditional software engineering---assumes capability is fixed and known. Give the system an input; receive a deterministic output.

This assumption fails catastrophically for large language models.

When an LLM receives a complex task, it does not possess fixed, queryable knowledge. Its effective capability on \emph{that specific task} depends on factors that emerge only through engagement: what context it accumulates, what connections it discovers, what it learns by exploring the problem space. An LLM asked to ``fix the authentication bug'' may have latent capability to solve it---but that capability becomes \emph{actual} only through investigation.

The execution-layer model ignores this. It assumes the AI either can or cannot perform the task, and proceeds accordingly. The result is a familiar failure mode:

\begin{itemize}
    \item AI reports low uncertainty (it ``thinks'' it knows)
    \item AI has low actual capability (it doesn't actually know)
    \item AI proceeds without investigation
    \item AI produces confident, wrong output
    \item User loses trust
\end{itemize}

We call this the \textbf{overconfidence trap}: the AI's self-assessed certainty exceeds its actual capability, leading to premature action without adequate exploration.

\subsection{Noetic Before Praxic}

This paper proposes a different model: \textbf{AI systems must investigate before acting}. We formalize this as two phases of AI cognition:

\begin{description}
    \item[Noetic phase:] Exploration, hypothesis formation, knowledge acquisition. The AI investigates the problem, reads relevant context, identifies unknowns, and builds genuine understanding.
    \item[Praxic phase:] Execution, implementation, action. The AI applies its acquired knowledge to produce output.
\end{description}

The critical insight is that these phases are not optional or interchangeable. \emph{Noetic must precede praxic}. An AI that acts without first investigating is an AI that acts on incomplete knowledge---and incomplete knowledge produces errors.

This is not a novel observation to practitioners who work with AI systems daily. It is, however, surprisingly absent from the research literature, which tends to evaluate AI on fixed benchmarks rather than dynamic, exploratory tasks.

\subsection{The Learning Delta}

If noetic precedes praxic---if investigation precedes action---then a measurable phenomenon should emerge: \textbf{AI capability should increase during task execution}.

An AI that begins a task with partial knowledge, investigates, and builds understanding will assess its own epistemic state differently at task-end than at task-start. The difference between these assessments---the \emph{learning delta}---measures genuine knowledge acquisition.

We present evidence from 479,686 observations across 638 sessions showing exactly this pattern. Across all 13 epistemic dimensions we track, AI systems show systematic confidence growth during task execution:

\begin{table}[h]
\centering
\begin{tabular}{lr}
\toprule
\textbf{Phase} & \textbf{Mean Capability Assessment} \\
\midrule
Pre-task (PREFLIGHT) & 0.663 \\
Post-task (POSTFLIGHT) & 0.863 \\
\textbf{Learning Delta} & \textbf{+0.200} \\
\bottomrule
\end{tabular}
\caption{Learning delta across 267 clean sessions.}
\label{tab:learning-delta-intro}
\end{table}

Filtering to 267 clean sessions (excluding default assessments and incomplete workflows), \textbf{93.6\% showed knowledge improvement}. This is not calibration error being corrected. This is learning being measured.

\subsection{Empirica: Epistemic Infrastructure for AI}

We introduce \textbf{Empirica}, a framework that operationalizes noetic-before-praxic through three mechanisms:

\begin{enumerate}
    \item \textbf{Epistemic Vectors}: A 13-dimensional representation of AI self-assessed cognitive state, covering knowledge (KNOW), uncertainty (UNCERTAINTY), engagement (ENGAGEMENT), and 10 additional dimensions derived from computational epistemology.

    \item \textbf{Epistemic Gates}: Decision points that evaluate whether the AI has sufficient knowledge to proceed. Gates enforce investigation when knowledge is low, preventing premature action.

    \item \textbf{The CASCADE Workflow}: A structured loop---PREFLIGHT $\rightarrow$ CHECK $\rightarrow$ POSTFLIGHT---that captures epistemic state before and after task execution, measures learning, and calibrates future assessments.
\end{enumerate}

Together, these mechanisms transform AI from an execution layer into a \textbf{learning system}---one that investigates before acting, measures its own knowledge acquisition, and improves through epistemic feedback.

\subsection{Contributions}

This paper makes the following contributions:

\begin{enumerate}
    \item \textbf{Conceptual}: We articulate the noetic-before-praxic principle and argue that treating AI as an execution layer is a category error that produces systematic failures.

    \item \textbf{Empirical}: We present 479,686 observations demonstrating that AI capability genuinely increases during task execution through investigation---the learning delta is real, not an artifact of assessment correction. 93.6\% of clean sessions show knowledge improvement.

    \item \textbf{Practical}: We introduce Empirica, an open-source framework implementing epistemic gates, self-assessment vectors, and calibration loops. The framework is production-tested across 638 sessions.

    \item \textbf{Theoretical}: We develop a computational epistemology for AI systems---a formal theory of how information-processing agents can assess and manage their own knowledge state.
\end{enumerate}

\subsection{Paper Structure}

\Cref{sec:related-work} reviews related work in AI self-assessment, metacognition, and calibration. \Cref{sec:theory} presents the theoretical framework: epistemic vectors, their grounding in transformer attention mechanisms, and the noetic-praxic distinction. \Cref{sec:implementation} describes the Empirica implementation. \Cref{sec:results} presents empirical results, including the learning delta analysis. \Cref{sec:sentinel} describes the Sentinel protocol for epistemic gating. \Cref{sec:discussion} discusses implications for AI deployment and alignment. \Cref{sec:conclusion} concludes.

\medskip
\noindent\emph{The core claim of this paper is simple: AI systems that investigate before acting outperform AI systems that act without investigating. The evidence is in the delta.}
